Dette projekt omhandler udvikling og implementering af et computerspil, der kører på en STM32-mikrocontroller. Spillet er skrevet i programmeringssproget C og er inspireret af det klassiske arkadespil "Space Invaders" med fokus på både funktionalitet og en optimal brugeroplevelse. Projektet benytter desuden en systemarkitektur, som gør mange af spillets komponenter, moduler og funktioner genanvendelige.

Spillet indeholder funktioner som meteorer med tyngdekraft, otte skudvinkler, power-ups samt point og liv vist på en LCD-skærm og kører hurtigt og effektivt. Spillet styres ved hjælp af det eksterne joystick 30010, hvor man både kan bevæge rumskibet, affyre skud og navigere rundt i menuen.

Projektet har givet værdifuld erfaring inden for C-programmering, samarbejde i udviklingsprocessen og softwarearkitektur. Denne rapport dokumenterer design, implementering, udfordringer og resultater og konkluderer, at spillet opfylder kravspecifikationerne, og leverer en unik og god brugeroplevelse.